\documentclass[a4paper,12pt,oneside,openany]{jsbook}
%
\usepackage{color}
\usepackage{amsmath,amssymb}
\usepackage{bm}
\usepackage[dvips,dvipdfmx]{graphicx}
\usepackage{verbatim}
\usepackage{wrapfig}
\usepackage{ascmac}
\usepackage{makeidx}
\usepackage{multirow}
\usepackage{enumerate}
\usepackage{url}
\usepackage{comment}
\usepackage{multirow}
\usepackage{array}

\usepackage[hang,small,bf]{caption}
\usepackage[subrefformat=parens]{subcaption}
\captionsetup{compatibility=false}


%%% 余白・文字数調整(左37mm, 右18mm, 上下共30mm, 文字数約40字/行, 行数約32行)
\setlength{\textwidth}{150truemm}      % テキスト幅: 210-(37+18)=155mm
\setlength{\fullwidth}{\textwidth}     % ページ全体の幅
\setlength{\oddsidemargin}{30truemm}   % 左余白
\addtolength{\oddsidemargin}{-1truein} % 左位置デフォルトから-1inch
\setlength{\topmargin}{30truemm}       % 上余白
\setlength{\textheight}{222truemm}     % テキスト高さ: 297-(30+30)=237mm
\addtolength{\topmargin}{-1truein}     % 上位置デフォルトから-1inch
\makeatother
%
\def\linesparpage#1{\baselineskip=\textheight
   \divide\baselineskip by #1}
\def\kcharparline#1{%
   \ifx\xkanjiskip\undefined%
   \jintercharskip 0mm plus 0.2mm minus 0.2mm
   \else
   \xkanjiskip 0mm plus 0.2mm minus 0.2mm
   \fi
   \settowidth{\textwidth}{}%
   \multiply\textwidth by #1}
%
\newcommand\figref[1]{\textbf{図~\ref{fig:#1}}}
\newcommand\tabref[1]{\textbf{表~\ref{tab:#1}}}
\newcommand{\bhline}[1]{\noalign{\hrule height #1}}
\renewcommand{\baselinestretch}{1.2}

\begin{document}

\begin{titlepage}

\begin{center}
{\Large 令和4年度　修士論文}\\ % タイトル
\vspace*{120truept}
{\LARGE 受験者固有層を有する深層学習モデルによる}\\ % タイトル
{\LARGE 複数記述式問題同時自動採点手法}\\ % タイトル
\vspace{100truept}

{\LARGE 電気通信大学　情報理工学研究科}\\ % 所属
\vspace{10truept}
{\LARGE 情報・ネットワーク工学専攻}\\ % 所属
\vspace{10truept}
{\LARGE 学籍番号 2031167}\\ % 学籍番号
\vspace{50truept}
{\LARGE HE　YIXIN}\\ % 著者
\vspace{50truept}
{\LARGE 指導教員：宇都雅輝　准教授}\\
\vspace{50truept}
{\LARGE 2022年7月29日}\\ % 提出日
\end{center}
\end{titlepage}
\kcharparline{40} % 一行を40字に

\frontmatter
%%% 論文要旨　（要旨は別ファイルで作成して統合する）
%\addcontentsline{toc}{chapter}{論文要旨}

近年，受験者の論理的思考力や表現力などの実践的な能力を評価する方法の一つとして，記述式問題が国内外の大規模試験を含む様々なテストで広く活用されている．一方で，特に大規模試験においては，採点の一貫性の確保が難しいことや採点に要する時間的・経済的コストが大きいことなどが記述式問題導入の課題として指摘されてきた．これらの問題を解決できる方法のひとつとして自動採点技術が近年注目されている．

記述式問題自動採点（Automated short-answer grading, ASAG）とは，短答記述式問題に対する回答文を人工知能技術を用いて自動的に採点する技術である．近年では，深層学習を用いた様々な自動採点モデルが提案され，高い精度を達成している．一般に深層学習自動採点モデルを含む従来の自動採点モデルは，記述式問題ごとに個別に収集されたデータセットを用いて訓練・構築され，問題ごとに独立に自動採点が行われる．しかし，現実のテストでは，同一テスト上で複数の記述式問題が出題されることがしばしばあり，そのような同一テスト上の問題群は受験者の特定の潜在的特性を測定するように設計されていると考えられる．そのような場合，複数の記述式問題におけるある受験者の得点は問題ごとに独立ではなく，受験者の特性という共通の要因に依存すると考えられる．したがって，複数の記述式問題の背後にある受験者固有の潜在特性を推定できれば，それは各問題に対する得点予測の有益な補助情報として機能すると期待できる．

そこで，本研究では，同一テスト上に同一の潜在特性を測定する複数の記述式問題が出題される場合を対象に，各受験者の複数の回答文からその受験者に固有の潜在特性を推定し，それを得点予測に活用する機構を持つ新たな深層学習自動採点モデルを提案する．提案モデルは，複数の短答式問題に対する回答文を同時に入力し，それらに対応する複数の得点を同時に出力する多入力多出力型の深層学習モデルとして定式化する．
また，本研究では，実データを用いた提案モデルの有効性評価実験から次のことを示した．
1）受験者固有の特徴量を加味することで得点予測の精度を改善できた．
2）BERTベースの提案モデルは，一定の訓練を受けた人間評価者による採点と同程度の精度を達成した．
3）提案モデルで抽出される受験者固有の特徴量は，数理モデルを用いたテスト理論の一つである項目反応理論に基づいて得点データから推定される受験者の能力値と相関しており，測定対象の受験者の能力を反映していると解釈できた．

%
%%% 目次
\tableofcontents
%
%
\mainmatter
%
%%% 本文ここから
\chapter{はじめに}\label{sec:intro}
近年，受験者の論理的思考力や表現力などの実践的な能力を評価する方法の一つとして，記述式問題が国内外の大規模試験を含む様々なテストで広く活用されている．一方で，特に大規模試験においては，採点の一貫性の確保が難しいことや採点に要する時間的・経済的コストが大きいことなどが記述式問題導入の課題として指摘されてきた．これらの問題を解決する方法のひとつとして，人工知能技術を用いて記述式回答に自動的に得点付けを行う記述式問題自動採点（Automated short-answer grading, ASAG）技術が注目されている．

従来の記述式問題自動採点手法は，専門家が事前に設計した特徴量を用いる特徴量ベースの手法と，深層学習モデルを用いる深層学習ベースの手法に大別される\cite{asag1,asag2,asag3,asag4,asag5}．

特徴量ベースの手法では，専門家が事前に設計した特徴量を説明変数とし，回帰モデルなどで得点を予測する．特徴量としては，総単語数や誤字脱字の数，文法エラーの数などが利用される．
特徴量ベースの手法は，古くから様々な手法が開発されてきており\cite{fearture1,fearture2,fearture7,fearture9}，代表的な手法としてはアメリカの教育試験機関ETS（Educational Testing Service）が開発したc-raterが知られている\cite{fearture9}．
一方で，一般に特徴量ベースの手法で高精度を実現するためには，採点対象の記述問題の特性に合わせた適切な特徴量を設計して使用することが必要となる．

このような特徴量の事前設計を回避して高精度を達成する手法として，近年では深層学習を用いた自動採点手法が多数提案されている\cite{cross-prompt1,cross-prompt2}．深層学習ベースの手法では，回答文の単語系列を入力として，得点予測に効果的な特徴量をモデル内部で自動的に構築する．深層学習ベースの手法はデータの特性に合わせた特徴量を自動的に獲得できるため，専門家による特徴量設計なしに高い精度の自動採点が実現できる．代表的な深層学習自動採点モデルとしては，畳み込みニューラルネットワークやLSTM（Long Short Term Memory）\cite{lstm1,lstm2}に基づくモデル，BERT（Bidirectional Encoder Representations from Transformers）\cite{BERT1,BERT2}に基づくモデルなどが知られており，それらの様々な拡張モデルも提案されている\cite{aes1,SKIPFLOWmodel,self-attention-basedmodel,DNN1,DNN2}．

上記のような従来の自動採点モデルは，記述式問題ごとに個別に収集されたデータセットを用いて訓練・構築され，問題ごとに独立に自動採点が行われる．しかし，現実のテストでは，同一テスト上で複数の記述式問題が出題されることがしばしばあり，そのような同一テスト上の問題群は受験者の特定の潜在的特性を測定するように設計されていると考えられる．
そのような場合，複数の記述式問題におけるある受験者の得点は問題ごとに独立ではなく，受験者の特性という共通の要因に依存すると考えられる．したがって，複数の記述式問題の背後にある受験者固有の潜在特性を推定できれば，それは各問題に対する得点予測の有益な補助情報として機能すると期待できる．

そこで，本研究では，同一テスト上に同一の潜在特性を測定する複数の記述式問題が出題される場合を対象に，各受験者の複数の回答文からその受験者に固有の潜在特性を推定し，それを得点予測に活用する機構を持つ新たな深層学習自動採点モデルを提案する．
提案モデルは，複数の短答式問題に対する回答文を同時に入力し，それらに対応する複数の得点を同時に出力する多入力多出力型の深層学習モデルとして定式化する．
具体的には，任意の既存の深層学習自動採点モデルを回答文のエンコーダーとして利用することで，各受験者の複数の回答文を分散表現ベクトルで表現し，それらを本研究で新たに提案するニューラルネットワーク層（受験者特徴抽出層と呼ぶ）に入力することで受験者固有の特徴量ベクトルを抽出する．
各問題に対する得点予測は，受験者固有の特徴量ベクトルと回答文の分散表現ベクトルを用いて行う．
提案モデルは従来の様々な自動採点モデルをエンコーダーとして利用できるが，本研究では代表的なモデルであるLSTMとBERTを用いたモデルを利用する．

また，本研究では，実データを用いた実験により提案モデルの有効性を評価する．実験の結果，次のことが明らかとなった．
1）受験者固有の特徴量を加味することで得点予測の精度を改善できた．
2）BERTベースの提案モデルは，一定の訓練を受けた人間評価者による採点と同程度の精度を達成した．
3）提案モデルで抽出される受験者固有の特徴量は，数理モデルを用いたテスト理論の一つである項目反応理論に基づいて得点データから推定される受験者の能力値と相関しており，測定対象の受験者の能力を反映していると解釈できた．

本論文の構成は以下の通りである．第2章では自動採点モデルについて説明する．第3章では提案モデルについて説明し，第4章では評価実験によって提案モデルの有効性を評価する．最後に第5章で本論文のまとめと今後の展望について述べる．


\chapter{自動採点モデル}

本章では，特徴量を用いた自動採点モデルと深層学習自動採点モデルについて紹介する．

\section{特徴量を用いた自動採点モデル}

特徴量ベースの自動採点手法は，古くから利用されてきた手法であり，専門家が事前に設計した特徴量（Hand-crafted features）を説明変数として，重回帰モデルや決定木，サポートベクターマシンなどの機械学習手法に基づいて得点を予測する.

最も代表的な特徴量ベースの短答記述式問題自動採点手法としては，アメリカの教育試験機関ETS（Educational Testing Service）が開発したc-raterが知られている\cite{fearture9}．
c-raterは全米学力調査とインディアナ州の学力調査2つの研究で使用されている．

また，Kleinら\cite{fearture10}は潜在意味解析技術 （Latent Semantic Analysis, LSA）モデルを用いた短答式問題自動採点手法を提案している．
この手法では，採点済み回答文と採点対象の回答文にLSAを適用し，採点対象文とトピック分布が類似した採点済み回答文を探索することで，自動採点を行う手法となっている．
これは，トピック分布を特徴量として利用した自動採点手法と解釈できる．

また，複数の文で構成される長めの文章を対象とした特徴量ベース自動採点モデルとしては，世界初の記述式自動採点システムPEG\cite{fearture3}やETSが開発したe-rater\cite{fearture7}，日本語の自動採点システムであるJESS\cite{fearture8}，近年の自動採点コンペティションで3位に入賞したEASE（Enhanced AI Scoring Engine）\cite{EASE}など，さまざまな手法が開発されてきた\cite{fearture1,fearture4,fearture5,fearture6}．


\section{深層学習を用いた自動採点モデル}\label{sec:lstm}

前節で紹介した特徴量ベースの手法において高精度を実現するためには，一般に，採点対象の記述問題の特性に合わせた適切な特徴量を設計して使用することが必要となる．
この問題を解決するアプローチの一つとして，回答文の単語系列を入力として，機械学習モデルを用いて得点予測に有効な特徴量を自動的に獲得するアプローチが提案されている．特に近年では，深層学習モデルを利用した自動採点手法が多数提案され，高精度を実現している\cite{lstm1,lstm2,lstm3,lstm4,lstm5,lstm6}．ここでは，代表的な深層学習自動採点モデルとしてLSTM（Long Short Term Memory）とBERT（Bidirectional Encoder Representations from Transformers）を用いたモデルを紹介する．その他の深層学習自動採点モデルについては付録で紹介する．


初期の深層学習自動採点モデルの一つであるLSTMを用いた自動採点モデル\cite{lstmAsgs}の概念図を図\ref{fig:LSTMmodel}に示す．
このモデルは，回答文の単語系列を入力として受け取り，多層のニューラルネットワークを通して得点の予測値を出力する．以下でモデルの各層について説明を行う．

\begin{figure}[t]
  \begin{center}
    \includegraphics[width=0.9\hsize]{Pic/pic2_LSTMmodel}
   \caption{LSTMを用いた自動採点モデルの構成}
 \label{fig:LSTMmodel}
  \end{center}
\end{figure}

\begin{enumerate}
\item {\bf 単語埋め込み層：}この層では回答文中の各単語を，類似した語彙が類似したベクトル表現として表現される「単語埋め込み表現」というベクトルに変換する．ここで，$\mathbf{w_{n}}$を回答文の$n\in\{1,2,...,N\}$番目の単語を表す$G$次元のone-hotベクトルとし，$\mathbf{E}$を$V\times G$次元の埋め込み行列とする．ただし，$N$は回答文中の単語数を表す．このとき，単語埋め込み層では，これらの内積$\mathbf{x_{n}}=\mathbf{E} \mathbf{w_{n}}$を計算することで各単語$\mathbf{w_{n}}$を$V$次元の埋め込みベクトル$\mathbf{x_{n}}$に変換する．

\item {\bf LSTM層：}この層ではLSTMに単語系列を入力することで，単語間の時系列的な関係を考慮したベクトルを求める．LSTMは，時系列的なデータ処理を行う深層学習アーキテクチャであるRNN（Recurrent Neural Network）の一種であり，従来のRNNの問題であった勾配消失問題（長い入力に対して指数関数的に勾配が小さくなってしまう問題）を解決するモデルの一つとして提案された．
LSTMでは長期的な文脈を保存する変数$\mathbf{c_{n}}$を用意し，その長期記憶変数$\mathbf{c_{n}}$に対して，古くなった記憶を削除し，新しい情報を新規追加する操作を順次行うことで時系列データの処理が行われる．
具体的には，入力ゲート，出力ゲート，忘却ゲートという３つのゲートを用いた以下の演算を繰り返して，長期記憶$\mathbf{c_{n}}$を更新していく．
\begin{equation}
\mathbf{i_{n}}= \sigma(\mathbf{W_{i}}\mathbf{x_{n}}+\mathbf{W^{\prime}_{i}}\mathbf{h_{n-1}}+\mathbf{b_{i}})
\end{equation}
\begin{equation}
\mathbf{f_{n}}= \sigma(\mathbf{W_{f}}\mathbf{x_{n}}+\mathbf{W^{\prime}_{f}}\mathbf{h_{n-1}}+\mathbf{b_{f}})
\end{equation}
\begin{equation}
\mathbf{\tilde{c}_n}= tanh(\mathbf{W_{c}}\mathbf{x_{n}}+\mathbf{W^{\prime}_{c}}\mathbf{h_{n-1}}+\mathbf{b_{c}})
\end{equation}
\begin{equation}
\mathbf{c_{n}}= \mathbf{i_{n}}\odot \mathbf{\tilde{c}_n}+\mathbf{f_{n}}\odot \mathbf{c_{n-1}}
\end{equation}
\begin{equation}
\mathbf{o_{n}}= \sigma(\mathbf{W_{o}}\mathbf{x_{n}}+\mathbf{W^{\prime}_{o}}\mathbf{h_{n-1}}+\mathbf{b_{o}})
\end{equation}
\begin{equation}
\mathbf{h_{n}}= \mathbf{o_{n}}\odot tanh(\mathbf{c_{n}})
\end{equation}

ここで，$\mathbf{W_{i}}$，$\mathbf{W_{f}}$，$\mathbf{W_{c}}$，$\mathbf{W_{o}}$は入力$\mathbf{x_{t}}$に対する重み行列，$\mathbf{W^{\prime}_{i}}$，$\mathbf{W^{\prime}_{f}}$，$\mathbf{W^{\prime}_{c}}$，$\mathbf{W^{\prime}_{o}}$は隠れ層$\mathbf{h_{n-1}}$に対する重み行列，$\mathbf{b_{i}}$，$\mathbf{b_{f}}$，$\mathbf{b_{c}}$，$\mathbf{b_{o}}$はそれぞれのバイアスベクトルである．また，$\odot$はアダマール積を表し，$\sigma(\cdot)$は式(2.7)に示すシグモイド関数，$tanh(\cdot)$は式(2.8)に示すハイパボリックタンジェントである．
\begin{equation}
\sigma(x)=\frac{1}{1+exp(-x)}
\end{equation}
\begin{equation}
tanh(x)=\frac{exp(x)-exp(-x)}{exp(x)+exp(-x)}
\end{equation}

\item {\bf プーリング層：}この層ではLSTM層の出力ベクトル系列${\mathbf{h_{1}},\mathbf{h_{1}},...,\mathbf{h_{N}}}$を時間方向に平均化する．この処理はMean over Time(MoT)と呼ばれ，具体的には次式で計算される．
\begin{equation}
\mathbf{M}=\frac{1}{N}\sum_{n=1}^{N}\mathbf{h_{n}}
\end{equation}

\item {\bf 回帰層：}この層では，プーリング層の出力ベクトルを重回帰に入力して得点に対応するスカラー値を求める．具体的には，次式で計算を行う．
\begin{equation}
s=\sigma(\mathbf{W}\mathbf{M}+\mathbf{b})
\end{equation}
ここで，$\mathbf{W}$，$\mathbf{b}$はそれぞれの重みとバイアスを表すパラメータである．
なお，この回帰層を通して得られる予測得点$s$は，シグモイド関数の利用により0から1の範囲の値となる．したがって，実際のデータの得点尺度はこれと異なる場合には，$s$を線形変換して実データの得点尺度に合わせる処理が必要となる．具体的には，1〜KのK段階得点の場合$Ks+1$と変換すればよい．
\end{enumerate}

\section{BERT}

本節ではBERTを用いた自動採点モデル\cite{BERT1,BERT4,Hybridmodel1,BERT6}を説明する．

BERTは，2018年にGoogle AI Languageチームが提案した事前学習済み深層学習モデルの一つであり\cite{BERT2}，質問応答（Question Answering）や自然言語推論（Multi Natural Language Inference）などの11種の自然言語処理タスクにおいて当時の最高性能を達成したモデルである\cite{BERT5}．
BERTは，RNNと同様に文字列のような時系列データを扱うモデルであるが，RNNとは異なりAttentionと呼ばれるメカニズムを利用して構成されている．具体的には，Attentionを多層に組み合わせた双方向Transformer\cite{BERT3}と呼ばれる機構で構成されている．

\subsection{Transformer}

TransformerモデルはVaswani et al.\cite{BERT3}によって提案されたモデルであり，RNNの代わりにSelf-AttentionとPosition Encodingと呼ばれる機構で設計されている．RNNは入力データを時系列順に読み込んで処理するため，時点$t$の入力データの処理が完了しないと時点$t+1$の処理ができず並列化が困難であった．一方で，Self-Attentionは，入力系列全体を一括で読み込んで系列全体の情報を同時に処理するため，並列計算が容易となることが知られている．さらに，Self-attentionを多層化・並列化した設計によって長期的かつ複雑な文脈を考慮できる利点も有する．そのような背景から，近年では，Transformerモデルに基づく様々な深層学習モデル（例えば，OpenAI GPT\cite{OpenAIGPT}やBERT，T5など）が提案されている．

Transformerモデルは，図\ref{fig:Transformer}のように，入力文章を符号化するエンコーダと出力文章を符号から文章に直すデコーダで構成される「エンコーダ・デコーダモデル」となっている．
なお，BERTはTransformerのエンコーダ部分のみを使用しているため，以下ではTransformerのエンコーダ部分を説明する．
\begin{figure}[t]
  \begin{center}
    \includegraphics[width=0.95\hsize]{Pic/pic_Transformer12}
   \caption{Transformerのエンコーダ・デコーダモデル構成}
 \label{fig:Transformer}
  \end{center}
\end{figure}

図\ref{fig:TransformerEncoder}に示したように，Transformerのエンコーダ部分は大きく二つの機構で構成されている．一つはMulti-head attentionであり，もう一つはFeedforward networkである．

\begin{figure}[t]
  \begin{center}
    \includegraphics[width=0.5\hsize]{Pic/pic_TransformerEncoder2}
   \caption{Transformerのエンコーダの構成}
 \label{fig:TransformerEncoder}
  \end{center}
\end{figure}

Multi-head attentionはself-attentionを拡張したアーキテクチャである．Self-attentionは入力行列$\mathbf{X}$に対して以下の演算を行う機構である．
\begin{equation}
\mathbf{Z}= \frac{\mathbf{Q}\mathbf{K_t}}{\sqrt{\mathbf{d_k}}}\mathbf{V}
\end{equation}
\begin{equation}
\mathbf{Q}= \mathbf{X}\mathbf{W_{Q}}
\end{equation}
\begin{equation}
\mathbf{K}= \mathbf{X}\mathbf{W_{K}}
\end{equation}
\begin{equation}
\mathbf{V}= \mathbf{X}\mathbf{W_{V}}
\end{equation}
ここで，入力行列$\mathbf{X}$は，$N$個の単語をそれぞれ512次元の埋め込みベクトルに変換して得られる$N\times512$であり，$\mathbf{W_Q}$,$\mathbf{W_K}$,$\mathbf{W_V}$は学習される$512\times64$次元の重み行列パラメータである．
なお，$\mathbf{d_k}$は$\mathbf{K}$の次元数である．

Multi-head attentionは，上記のSelf-attentionを複数並列に用意して結合した機構である．例えば，Self-attentionを8個用意した場合，Multi-head attentionは以下の式で計算される．
\begin{equation}
\mathbf{Z}= Concatnate(\mathbf{Z_1},\mathbf{Z_2},\cdots,\mathbf{Z_8})\mathbf{W_o}
\end{equation}
ここで，$\mathbf{W_o}$は線形変換のための重みパラメータ行列であり，$Concatnate$はベクトルの連結を表している．

Multi-head attentionで得られた行列$\mathbf{Z}$は，全結合層で構成されるFeedforward networkに入力して，その出力行列が次の層のTransformerブロックに入力される．
Transformerのエンコーダーは，以上の処理を多層に積み上げたモデルとして設計されている．

\subsection{BERTモデルの事前学習とファインチューニング}

BERTは，上記のように設計されたTransformerモデルのエンコーダーを，大量の文章データで事前学習（pre-training）したモデルである．具体的には，WikipediaとBookCopusから集めた10億語以上の大規模な文章データを利用して事前学習されている．事前学習には，「Masked Language Model」と「Next Sentence Prediction」の二つの教師なし学習タスクが採用されている．
ここで，Masked Language Modelとは，文字列の一部の単語を「Mask」トークンに置き換え，マスクされている単語を予測するタスクであり，
Next Sentence Predictionとは入力文字列のペアを受け取り，それらが連続した文であるかを予測するタスクである．

BERTを自動採点を含む別の言語処理タスクに利用するためには，ファインチューニングを行う必要がある．ファインチューニングとは，事前学習で得られたパラメータを初期値として，対象とする言語処理タスクに関するラベル付き学習データを用いてモデルを再学習をすることを意味する．

\begin{figure}[t]
  \begin{center}
    \includegraphics[width=0.5\hsize]{Pic/pic7_BERT}
   \caption{BERTを用いた自動採点モデルの構成}
 \label{fig:BERT}
  \end{center}
\end{figure}

自動採点タスクにBERTを適用するための手順は以下の通りである：
\begin{enumerate}
\item 各回答文の先頭に特殊タグ（[CLS]）を追加する．
\item 図\ref{fig:BERT}のように，[CLS]タグに対するBERTの出力を用いて予測得点を計算する回帰層を追加する．
\item このように出力層を追加したBERTモデルを，得点付き回答文データを用いてファインチューニングする．
\end{enumerate}

これまで，BERTは短答記述式問題や小論文問題などの文章自動採点に広く利用され，高精度を達成している\cite{DNN1,DNN2,DNN3,DNN4,DNN6}．

\section{深層学習自動採点モデルの学習}\label{sec:train}

前述の深層学習ベースのモデルを利用するためには，大量の得点付き回答文のデータセットを用いてモデルの学習やファインチューニングを行う必要がある．
具体的には，モデルの予測得点と真の得点の誤差を表す以下の平均二乗誤差（Mean Square Error: MSE）を最小化するように学習を行うことが一般的である．
\begin{equation}
MSE(\bm{y},\hat{\bm{y}})=\frac{1}{I}\sum_{i=1}^{I}(y_i-\hat{y}_i)^2
\end{equation}
ここで，$I$はトレーニングデータセット中の回答文の数を表し，$y_{i}$ と $\hat{y}_{i}$は$i$番目の回答文に対する真の得点と予測得点を表す．

なお，\ref{sec:lstm}節で述べたように，従来モデルの多くは回帰層にシグモイド関数を用いているため，予測される得点 $\hat{y}_{i}$ は $[0, 1]$ の範囲で与えられる．
このような場合，真のスコア$y_{i}$を範囲$[0, 1]$に正規化した後でモデル学習を行う必要がある．

一度モデルの訓練が完了すれば，新しい回答文を自動採点することができる．具体的には，訓練されたモデルに新たな回答文のデータを入力することで予測得点を得ることができる．ただし，上述の通り，予測値は[0,1]で与えられるため，ここでも元の得点尺度への線形変換が必要となることに注意する．

\section{従来手法の問題点}\label{sec:old}

\begin{figure}[t]
  \begin{center}
    \includegraphics[width=0.95\hsize]{Pic/pic3_3Lstm}
   \caption{複数問題に対する従来モデルの構成}
 \label{fig:3Lstm}
  \end{center}
\end{figure}

\ref{sec:intro}章で説明したように，従来手法では，複数の記述式問題を自動採点する場合，記述式問題ごとに個別に収集されたデータセットを用いてモデルを訓練・構築する．
図\ref{fig:3Lstm}に，3つの記述式問題がある場合の従来手法の適用法の概念図を示す．このように，従来手法では，3つの問題に対してそれぞれ独立に自動採点モデルを構築する．

一方で，現実のテストでは，図\ref{fig:3qusetions}のように同一テスト上で複数の記述式問題が出題されることがしばしばあり，そのような同一テスト上の問題群は受験者の特定の潜在的特性を測定するように設計されていると考えられる．
そのような場合，複数の記述式問題におけるある受験者の得点は問題ごとに独立ではなく，受験者の特性という共通の要因に依存すると考えられる．したがって，複数の記述式問題の背後にある受験者固有の潜在特性を推定できれば，それは各問題に対する得点予測の有益な補助情報として機能すると期待できる．
本研究の目標は，複数の記述式問題の背後にある受験者固有の潜在特性を，その受験者の複数の回答文から抽出する機構を開発し，その有効性を検討することである．

\begin{figure}[t]
  \begin{center}
    \includegraphics[width=0.4\hsize]{Pic/pic1_3qusetions}
   \caption{複数の短答記述式問題で構成されるテスト}
 \label{fig:3qusetions}
  \end{center}
\end{figure}

\chapter{提案モデル}

\ref{sec:old}節で述べた既存手法の問題点を解決するために，本研究では，同一テスト上に同一の潜在特性を測定する複数の記述式問題が出題される場合を対象として，各受験者の複数の回答文からその受験者に固有の潜在特性を推定し，それを得点予測に活用する機構を持つ新たな深層学習自動採点モデルを提案する．
提案モデルは，複数の短答式問題に対する回答文を同時に入力し，それらに対応する複数の得点を同時に出力する多入力多出力型の深層学習モデルとして定式化する．

\begin{figure}[t]
  \begin{center}
    \includegraphics[width=0.9\hsize]{Pic/pic8_proModel}
   \caption{LSTMベースの提案モデルの構成}
 \label{fig:ModelLSTM}
  \end{center}
\end{figure}

\begin{figure}[t]
  \begin{center}
    \includegraphics[width=0.9\hsize]{Pic/pic10_proModelBERT}
   \caption{BERTベースの提案モデルの構成}
 \label{fig:ModelBERT}
  \end{center}
\end{figure}

提案モデルは，問題固有エンコーダー，受験者特徴抽出層，問題固有出力層の3つの部分から構成されている．
各要素の処理は以下の通りである．
\begin{enumerate}
\item {\bf 問題固有エンコーダー：} ここでは，各回答文を従来の自動採点モデルと同様の構造を持つ深層学習モデルに入力し，固定長の実数値ベクトルで表現される文章レベルの分散表現ベクトルを生成する．
問題固有エンコーダーには，様々な自動採点モデルを利用できる．例として，問題固有エンコーダーにLSTMベースのモデルを用いた場合の構造を図\ref{fig:ModelLSTM}に示す．
このエンコーダーでは，従来のLSTMベースの自動採点モデルと同様に，入力文を1）単語埋め込み層，2）LSTM層，3）プーリング層の順に処理して，分散表現ベクトルを求める．
また，図\ref{fig:ModelBERT}に，BERTベースのモデルを問題固有エンコーダーとして用いた場合の構造を示す．
この場合には，[CLS]タグ付きの回答文をBERTベースのモデルに入力し，[CLS]タグに対応する出力ベクトルを分散表現ベクトルとして抽出する．
なお，いずれの場合もエンコーダは問題ごとに固有であり，回答文は問題ごとに異なるエンコーダで処理される．

\item {\bf 受験者特徴抽出層：} この層では，問題固有エンコーダーで得られた複数の回答文に対応する複数の分散表現ベクトルを用いて，受験者固有の潜在的特性の抽出を行う．
ここで，採点対象の記述式問題数を$Q$とし，問題固有エンコーダーで得られた$Q$個の回答文に対応する分散表現ベクトルをそれぞれ$\bm{F}_1, \bm{F}_2, \ldots, \bm{F}_Q$とする．
このとき，受験者特徴抽出層では，これらの分散表現ベクトル$\bm{F}_1, \bm{F}_2, \ldots, \bm{F}_Q$を連結したベクトル$\bm{F}$を作成し，それを以下の全結合層に通して低次元のベクトル$\bm{F}'$に変換する．
\begin{equation}
\bm{F}'=\sigma(\mathbf{W}^{(e)}\bm{F}+\mathbf{b}^{(e)})
\end{equation}
提案モデルでは，このベクトル$\bm{F}'$を，受験者固有の潜在的特性を表すベクトルとみなす．
受験者特性ベクトル$\bm{F}'$は各回答文の分散表現ベクトル$\bm{F}_q$と以下のように連結され，次の問題固有出力層に渡される．
\begin{equation}
\bm{F^{(c)}_q} = \bm{F}' \oplus \mathbf{F_{q}}
\end{equation}
ここで，$\bm{F}'$と$\bm{F}_q$を連結する処理は，受験者固有の特性を抽出するために重要であることに注意してほしい．
受験者特徴抽出層から出力層に渡されるベクトルを$\bm{F}'$のみとした場合，$\bm{F}'$には受験者固有の特性と個別の回答文に固有の特徴量が同時に含まれることになり，抽出を目指した受験者固有の特性に焦点化できない．
一方で，$\bm{F}'$と$\bm{F}_q$を連結して入力すれば，問題固有の特徴量は$\bm{F}_q$に含まれるため，$\bm{F}'$には複数の回答文の背後にある共通の特徴量のみを抽出するように動作すると期待できる．
このような構成は，深層学習分野において残差接続（Residual Connection）と呼ばれ，広く利用されている．

\item {\bf 問題固有出力層：} ここでは，問題ごとの回帰層を用いて，受験者特徴抽出層から得られた各ベクトル$\bm{F^{(c)}_q}$をスカラー値の得点$s$に変換する．
具体的には，従来の自動採点モデルと同様に，シグモイド関数を活性化関数に用いて以下で計算する．
\begin{equation}
s=\sigma(\mathbf{W}\bm{F^{(c)}_q}+\mathbf{b})
\end{equation}
\end{enumerate}

提案モデルの学習には誤差逆伝播法を使用し，損失関数としては以下の$MSE$関数を用いる．
\begin{eqnarray} 
Loss =\frac{1}{Q\cdot I}\sum_{q=1}^Q \sum_{i=1}^I (y_{iq}-\hat{y}_{iq})^2, 
\end{eqnarray}
ここで、$I$はトレーニングデータ中の受験者数，$y_{iq}$と$\hat{y}_{iq}$は問題$q$に対する$i$番目の受験者の回答文の真の得点と自動採点モデルの予測得点を示す．
なお，\ref{sec:train}節でも説明したように，回帰層にシグモイド関数を用いているため，モデル学習を行う前に真得点は$[0, 1]$尺度に正規化する必要がある．
また，予測を行うときには，予測された得点$s$を元の得点範囲に線形変換する必要がある．

\chapter{実験}

本章では，実データを用いた実験により，提案モデルの有効性を評価する．

\section{実データ}\label{sec:actdata}

本実験では，ベネッセ教育総合研究所が日本人大学生を対象に実施した読解テストの回答データを使用した．この読解テストは学習到達度調査（PISA）における「読解力」の定義に基づいてベネッセ教育総合研究所が独自開発したものである．このテストでは，グラフやイラストを含む200字〜2000字程度のテキスト素材が与えられ，そのテキストから適切に情報を抽出・推論・理解できるかを問う記述式問題が出題される．
このデータセットは3つの短答式問題に対する$511$の日本人大学生受験者の回答で構成されている．
採点はまず事前に訓練された2名の評価者によって行われ，その後，3人目の評価者が2名の採点結果を参考に最終的な採点結果を決定することで行われた．
得点は，質問1が4段階評価，質問2と3が5段階評価で与えられている．
表\ref{table:kiso}に各問題の回答文の文字数の平均と標準偏差を示す．
また，表\ref{table:kiso_score}に，各短答記述式問題に対する得点の平均と標準偏差，および各得点カテゴリの出現頻度を示した．

\begin{table*}[t]
\begin{center}
\caption{回答文の文字数に関する基礎統計量と得点の段階数}
\setlength{\tabcolsep}{5.pt}
\begin{tabular}{ccccc}  
\bhline{1pt}
 &  & 問題1& 問題2& 問題3\\
\hline
\multirow{2}*{文字数}& 平均& 27.1& 33.0& 49.6\\
\cline{2-5}
 & 標準偏差& 7.51& 10.60& 13.60\\
\hline
\multicolumn{2}{c}{\multirow{2}*{得点の段階数}}& 4& 5& 5\\
\multicolumn{2}{c}{~}& (0,1,2,3)& (0,1,2,3,4)& (0,1,2,3,4)\\
\bhline{1pt}
\end{tabular}
\label{table:kiso}
\end{center}
\end{table*}

\begin{table}[t]
\tabcolsep = 1.5mm
\center
\caption{得点データに関する基礎統計量}
\begin{tabular}{llcccccccccccc}
\bhline{1pt}
 &  &  & \multicolumn{5}{c}{得点カテゴリの出現頻度} \\
 \cline{4-8} 
 &  $平均$ & $標準偏差$ & 1 & 2 & 3 & 4 & 5 \\
\bhline{1pt}
問題1& 2.307 & 0.813 & 82 & 222 & 175 & 32 & -  \\
問題2& 3.223 & 1.207 & 63 & 44 & 214 & 96 & 94 \\
問題3& 3.838 & 1.376 & 35 & 65 & 126 & 7 & 278 \\
\bhline{1pt}
\end{tabular}
\label{table:kiso_score}
\end{table}

\section{得点データの因子分析}

提案モデルでは，複数の記述式問題が共通の受験者特性（能力など）を測定していることを仮定している．
そこで，今回のデータがこの仮定を満たすかどうかを確認するために，得点データに因子分析を適用してデータの背後にある潜在尺度の次元数を検証する．
因子分析に基づく潜在尺度の次元数の検証には，スクリープロットを用いて固有値が1を下回る直前の次元数を最適値とするカイザー基準で判定することが多い．

本実験で使用するデータから得られたスクリープロットを図\ref{fig:Rplot}に示す．
図の横軸は潜在尺度の次元数，縦軸は固有値を表す．
図から，次元数を1とした場合には固有値が1を超えており，次元数が2以上となると固有値が１を下回っていることから，本実験で用いた複数の記述式問題は一次元の共通する受験者特性を測定していると解釈できる．

\begin{figure}[t]
  \begin{center}
    \includegraphics[width=0.4\hsize]{Pic/picData_Rplot}
   \caption{得点データから得られたスクリープロット}
 \label{fig:Rplot}
  \end{center}
\end{figure}

\section{従来モデルと提案モデルの比較}

ここでは，実データを用いて従来モデルと提案モデルにおける自動採点の精度を評価する．

\subsection{実験手順}

比較するモデルは，1）LSTMベースの従来モデル，2）LSTMベースの提案モデル，3）BERTベースの従来モデル，4）BERTベースの提案モデル，とした．
実験は5分割交差検証法で行った．具体的には，60\%のデータを訓練データ，20\%のデータをハイパーパラメータチューニングに用いる検証データ，残りの20\%のデータを精度評価に利用する評価データとし，訓練データ/検証データ/評価データの選び方を変えながら5回評価を行い，精度指標値の平均を求めた．
評価指標には，自動採点の研究で一般的に用いられる二乗平均平方根誤差（RMSE）とピアソン相関係数，重み付きカッパ係数（Quadratic Weighted Kappa: QWK）を使用した．

LSTMベースモデルはPython-Kerasで実装した．なお，オーバーフィッティングを防ぐため，ドロップアウト正則化を用いた．具体的には，単語埋め込み層の出力ドロップアウトに$0.5$を，LSTMのリカレントドロップアウトに$0.2$を設定した．
学習アルゴリズムにはAdamを用い，ミニバッチサイズを$64$，最大エポック数を$100$とした．
ハイパーパラメータは検証データを用いたグリッドサーチにより最適化した．
グリッドサーチにおける候補値としては，単語埋め込み次元数を$50$と$100$，LSTMの次元数を$100$と$200$，受験者固有特性の次元数を$10$と$20$とした．
グリッドサーチの結果，LSTMベースの従来モデルでは単語埋め込み次元数$50$とLSTM次元数$200$が最適値となり，LSTMベースの提案モデルでは単語埋め込み次元数$100$，LSTM次元数$100$と受験者固有特性の次元数$20$が最適値となった．
BERTベースのモデルは，PytorchとTransformersライブラリで実装した．
事前学習BERTモデルには，日本語テキストを事前学習させた{\it base}サイズモデル (cl-tohoku/bert-base-japanese-v2\footnote{https://huggingface.co/cl-tohoku/bert-base-japanese-v2})を使用した．ファインチューニングは，ミニバッチサイズを$8$，エポック数を$200$として，AdamWアルゴリズムで行った．
受験者固有特性の次元数は，候補値を$10$と$20$としてグリッドサーチで検証し，$20$次元を最適値として決定した．
表\ref{table:hyperparameter}に，各モデルの最適なハイパーパラメータをまとめた．以降の実験では，これらの値を利用した評価結果について述べる．

\begin{table}[t]
\tabcolsep = 1.5mm
\center
\caption{グリッドサーチで決定した各モデルのハイパーパラメータの設定}
\begin{tabular}{ccccc}
\bhline{1pt}
 &  単語埋め込み層&  LSTM層& 受験者固有層\\
\bhline{1pt}
LSTMベース従来モデル&  50& 200& -\\ 
LSTMベース提案モデル&  100& 100& 20\\ 
BERTベース提案モデル&  -& -& 20\\ 
\bhline{1pt}
\end{tabular}
\label{table:hyperparameter}
\end{table}

\subsection{実験結果}

表~\ref{tb:result}に各モデルにおける自動採点の精度を示す．

まず，LSTMベースとBERTベースのモデルを比較すると，全ての場合でBERTベースのモデルのRMSEが低く，相関とQWKが高いことがわかる．
これは，LSTMベースのモデルがBERTベースのモデルに比べて高い自動採点精度を達成していることを意味する．
この傾向は，他の多くの研究の結果とも一致している~\cite{asag6,asag7,asag8,asag9,okano}. 

次に，提案モデルと従来モデルを比較すると，全てのケースで提案モデルが高い精度を示していることがわかる．
ここで，提案モデルによる改善量が統計的に有意であるかを確認するために，提案モデルと従来モデルの予想精度の平均値について対応のある\textit{t}検定を実施した．
表~\ref{tb:result}の$p$列に\textit{t}検定で得られた$p$値を示した．
この結果から，全ての場合において，提案モデルが有意水準5\%で従来モデルより高精度であったことが確認できた．

なお，表~\ref{tb:result}から，BERTベースモデルを使用した場合には，従来モデルと提案モデルで改善量がやや小さいことが読み取れるが，
これは，次の節で示すようにBERTベースのモデルの精度が人間による採点の精度に十分に近づいており，大幅な精度改善が難しくなっている可能性を示唆している．

\begin{table*}[t]
\tabcolsep = 1mm
\center
\caption{従来モデルと提案モデルの実験結果}
\label{tb:result}
\begin{tabular}{ll ccc cc}
\bhline{1pt}
 &  & 問題1 & 問題2 & 問題3 & $平均$ & $p$ \\
\bhline{1pt}
RMSE & LSTMベース従来モデル & 0.654 & 0.720 & 1.000 & 0.791 &  \multirow{2}{*}{0.026} \\
 & LSTMベース提案モデル &\bf 0.633 &\bf 0.673 &\bf 0.949 &\bf 0.752 & \\
 \cline{2-7}
 & BERTベース従来モデル & 0.383 & 0.462 & 0.772 & 0.539 & \multirow{2}{*}{0.011}  \\
 & BERTベース提案モデル &\bf 0.366 &\bf 0.451 &\bf 0.753 &\bf 0.523  \\
\hline
相関係数 & LSTMベース従来モデル & 0.616 & 0.820 & 0.715 & 0.717 & \multirow{2}{*}{0.032}  \\
 & LSTMベース提案モデル &\bf 0.660 &\bf 0.841 &\bf 0.738 &\bf 0.746  \\
 \cline{2-7}
 & BERTベース提従来モデル & 0.884 & 0.930 & 0.838 & 0.884 & \multirow{2}{*}{0.031} \\
 & BERTベース提案モデル &\bf 0.892 &\bf 0.933 &\bf 0.843 &\bf 0.889 \\
\hline
QWK & LSTMベース従来モデル & 0.547 & 0.774 & 0.665 & 0.662 &  \multirow{2}{*}{0.008} \\
 & LSTMベース提案モデル & 0.612 & 0.826 & 0.714 & 0.717 & \\
 \cline{2-7}
 & BERTベース従来モデル & 0.858 & 0.891 & 0.795 & 0.848 & \multirow{2}{*}{0.033}  \\
 & BERTベース提案モデル &\bf 0.880 &\bf 0.907 &\bf 0.834 &\bf 0.874  \\
\bhline{1pt}
\end{tabular}
\end{table*}

\section{提案モデルと人間評価者の比較}


本節では，人間の評価者による採点の精度を検証を通して，提案した自動採点モデルの精度について考察を行う．
\ref{sec:actdata}節で説明したように，本実験で使用した得点は，訓練された評価者2名が付けた得点に基づいて別の1名の評価者が最終的に調整を行うことで作成されている．
そこで，ここでは，2名の訓練された評価者の得点と，最終的に調整された得点（真得点と呼ぶ）との関係を調べた．
指標には，前節の精度評価実験と同様に，RMSEと相関係数，QWKを用いた．

結果を表~\ref{tb:result2}に示す．
表の「平均」の行は，2名の評価者の精度の平均を示している．

この表から，2名の評価者では，それぞれに採点の傾向が異なることが読み取れる．
例えば，評価者~1は問題~3については高精度に採点する傾向があるが，問題~1に対しては誤差が大きい傾向がある．
一方で，評価者~2は問題~1については的確に評価できているが，問題~3については誤差が大きい．

次に，「平均」行に示された結果を表~\ref{tb:result}で示した自動採点モデルによる精度と比較すると，
BERTベースのモデルの精度が2名の人間評価者の平均的な精度と概ね同等の値を示していることがわかる．
このことから，少なくとも本実験で使用したデータにおいては，BERTベースのモデルはすでに人間の評価者と同等の水準を達成しており，精度の大幅な向上は困難であると解釈できる．
ただし，このような状況においても，提案モデルは従来モデルと比べて有意な精度の向上に成功しており，数値上の改善はやや小さく見えるものの提案モデルは精度改善に有効であると結論づけられる．

\begin{table*}[t]
\center
\caption{真得点と人間評価者得点の関係性} \label{tb:result2}
\begin{tabular}{c l cccc}
\bhline{1pt}
& & 問題1 & 問題2 & 問題3 & $平均$ \\
\bhline{1pt}
RMSE& 評価者1& 0.709 & 0.380 & 0.543 & 0.544 \\
& 評価者2 & 0.286 & 0.364 & 0.821 & 0.491 \\
\cline{2-6}
& $平均$ & 0.498 & 0.372 & 0.682 & 0.517 \\
\hline
相関係数& 評価者1& 0.839 & 0.951 & 0.925 & 0.905 \\
& 評価者2& 0.972 & 0.957 & 0.834 & 0.921 \\
\cline{2-6}
& $平均$ & 0.905 & 0.954 & 0.880 & 0.913 \\
\hline
QWK& 評価者1& 0.818 & 0.951 & 0.916 & 0.895 \\
& 評価者2& 0.971 & 0.955 & 0.815 & 0.914 \\
\cline{2-6}
& $平均$ & 0.894 & 0.953 & 0.865 & 0.904 \\
\bhline{1pt}
\end{tabular}
\end{table*}

\section{受験者固有特性の分析}

本節では，提案モデルの受験者特徴抽出層から得られた受験者固有特性について考察を行う．
これまでに説明したように，受験者固有特性は，複数の問題に共通する受験者の潜在的な特性を反映することが期待される．
そこで，本節では，数理モデルを用いたテスト理論の一つとして知られる項目反応理論（IRT: Item Response Thepry）~\cite{IRT1}を用いて得点データから各受験者の潜在特性を推定し，その推定値と提案モデルで抽出した受験者固有特性との関係を分析する．
本研究で扱うような多値のリッカート型データに適用できる項目反応理論のモデルとしては，評定尺度モデル（Rating Scale Model：RSM）~\cite{IRT3}や部分採点モデル（Partial Credit Model：PCM）~\cite{IRT4}，一般化部分採点モデル（Generalized Partial Credit Model：GPCM）~\cite{IRT2}などが知られている．
ここでは，RSMやPCMの一般化モデルとみなせるGPCMを用いた．
GPCMは，受験者$j$が問題$i$において得点$k\in\{1\cdots K\}$を得る確率$P_{ijk}$を次式で定義する．
\begin{eqnarray} 
P_{ijk} = \frac{\exp{\sum_{m=1}^k[D\alpha_i(\theta_j-\beta_{i}-d_{im})]}}{\sum_{l=1}^K\exp{\sum_{m=1}^l[D\alpha_i(\theta_j-\beta_{i}-d_{im})]}}
\end{eqnarray}
ここで，$\theta_j$は受験者$j$の能力を，$\alpha_i$は問題$i$の識別力を表す．$\beta_{i}$は問題$i$の困難度を表す位置パラメータであり，$d_{im}$は問題$i$において得点$k$を得る困難度を表すステップパラメータである．なお，モデルの識別性のために，$d_{i1} = 0,\sum_{k=2}^K d_{ik} = 0:\forall i$を所与とする．
これらの項目パラメータと能力パラメータはテストへの反応データから推定することができる．
本実験では，GPCMのパラメータ推定をマルコフ連鎖モンテカルロ（Markov Chain Monte Carlo: MCMC）法を用いた期待事後確率推定（Expected A Posteriori:EAP）で求めた．

\begin{table*}[t]
\center
\caption{GPCMで推定された受験者の1次元潜在特性値とBERTベースの提案モデルで得られた20次元の受験者固有特性との相関} \label{tb:relation}
\begin{tabular}{ccc|ccc}
\hline
次元& 相関& & 次元& 相関&  \\
\hline
1& -0.440& **& 11& 0.159& ** \\
2& -0.289& **& 12& 0.125& ** \\
3& -0.001& & 13& -0.146& ** \\
4& -0.010& & 14& -0.311& ** \\
5& 0.031& & 15& 0.104& * \\
6& 0.101& *& 16& -0.106& * \\
7& -0.266& **& 17& -0.223& ** \\
8& 0.074& & 18& -0.240& ** \\
9& 0.057& & 19& -0.248& ** \\
10& 0.040& & 20& 0.168& ** \\
\hline
\end{tabular}
\end{table*}

表~\ref{tb:relation}には，GPCMで推定された受験者の1次元潜在特性値と，BERTベースの提案モデルで得られた20次元の受験者固有特性の各次元との間で計算した相関係数を表す．
表~\ref{tb:relation}において，**は相関係数の有意差検定の$p$値が0.01未満で，*は0.05未満であったことを示す．

表~\ref{tb:relation}より，提案モデルで得られた受験者固有特性の多くの次元がGPCMで推定された潜在特性値と有意な相関を示していることがわかる．
このことから，提案モデルでは，複数の記述式問題に対する得点データの背後にある共通的な受験者特性を，その受験者の回答文だけから推定できていることがわかる．
また，表~\ref{tb:relation}の結果は，得点データから項目反応理論で推定される潜在特性（一般に，受験者の能力と解釈される）とは相関しない次元もいくつかあることがわかる．
これらの次元は，得点データだけではわからない受験者固有の特性（語彙選択の傾向や句読点などを含む文章の傾向の差異など）を反映している可能性がある．
その詳細な分析については今後の課題としたい．

\chapter{まとめ}

本研究では，同一テスト上に同一の潜在特性を測定する複数の記述式問題が出題される場合を想定し，各受験者の複数の回答文からその受験者に固有の潜在特性を推定し，それを得点予測に活用する機構を持つ新たな深層学習自動採点モデルを提案した．
具体的には，各受験者の複数の回答文を任意の既存の深層学習自動採点モデルに入力して分散表現ベクトルに変換し，それらを入力として受験者固有の特徴量ベクトルを抽出する受験者特徴抽出層を開発した．
本研究では，分散表現ベクトルを構築するエンコーダーとして従来のLSTMベースとBERTベースの自動採点モデルを利用した提案モデルを実装し，実データ実験を通してその有効性を評価した．
実験の結果，提案モデルを用いることで自動採点の精度が向上することが確認できた．
また，BERTベースの提案モデルは，人間の評価者と同程度の採点精度を達成することが確認された．
また，提案モデルで複数の回答文から計算される受験者固有特性は，項目反応理論を用いて得点データから推定される受験者の潜在特性と強く相関することがわかり，複数の記述式問題が測定している潜在特性を適切に反映していることが確認できた．

なお，4.5節で示したように，提案モデルで抽出される受験者固有特性は，得点データだけではわからない受験者固有の特性（語彙選択の傾向や句読点などを含む文章の傾向の差異など）を反映している可能性も示唆された．
今後は，それらの特性の詳細な分析を行っていきたい．
また，提案モデルの問題固有エンコーダーとして，本研究では，LSTMベースの自動採点モデルとBERTベースの自動採点モデルを用いたが，付録にも示したようにこの他にも様々な自動採点モデルが提案されている．
今後は，提案モデルを様々な深層学習ベースの自動採点モデルに適用したい．
さらに，今後はより多様なデータセットを用いた実験で提案モデルの性能を評価していきたい．

%
%
%%% 付録
%\appendix
\chapter*{付録：様々な深層学習自動採点モデル}

本研究では，LSTMベースモデルとBERTベースモデルに基づいた複数問題の同時自動採点モデルを検討したが，近年ではこれら以外にも様々な深層学習自動採点モデルが提案されている．
ここでは，代表的なモデルをいくつか紹介する．

\subsection*{短答記述式問題のための深層学習自動採点モデル}\label{sec:other_asag}

\subsubsection*{Semantic feature-wise transformation relation network}

一般に，短答記述式問題自動採点のデータセットは，問題文と模範解答（またはルーブリック），受験者の採点済み回答文で構成されている．
一般的な自動採点モデルは，受験者の採点済み回答文のみを用いて，モデルの訓練や得点の予測を行うが，
Liら\cite{ASAG_SFRN}は問題文や模範解答，ルーブリックも利用して自動採点を行うモデルを提案している．
このモデルは，問題文と参照答案またはルーブリック，受験者の回答文について類似性を求め，それに基づいて得点を予測するように設計されている．

\subsubsection*{Geometric average normalized-longest common subsequence}

上述のSemantic feature-wise transformation relation networkのように，模範解答と受験者の回答文の類似性を用いて採点する短答式問題自動採点手法が提案されてきた．
しかし，この方法を利用するためには，受験者の回答文の多様性に対応できるような多様な模範解答と，それらの模範解答と受験者の回答文の類似性を測る適切な指標の設計が必要となる．
そこで，文献\cite{asag11}では，文章自動要約技術の一つである最大限界関連度法を用いて，事前に用意された模範解答と採点対象となる受験者の回答文から，多様なパタンの模範解答を生成する手法を提案している．
さらに，模範解答と受験者の回答文の類似性指標としては，幾何平均正規化-最長共通部分列と呼ばれる類似度指標を提案している．

\subsubsection*{Mixed-Formatテストへの適用を想定した短答記述式問題自動採点モデル}

Uto \& Uchida\cite{uchida}\cite{uchida2}は短答式記述式問題が一般に客観式問題を含むテストの一部として出題される場合を想定し，項目反応理論\cite{IRT1}を用いて客観式問題から求められる受験者の能力推定値を組み込んだ新たな深層学習自動採点モデルを提案している．具体的には，深層学習自動採点モデルの内部で獲得される回答文の分散表現ベクトルに，客観式問題への正誤データから項目反応理論を用いて推定される受験者の能力値を統合して，回答文の得点を予測するモデルを開発している．

\subsubsection*{短答記述式問題自動採点のためのデータ拡張}

短答記述式問題自動採点のデータ拡張を行う方法も提案されている．データ拡張とは，既存のデータセットから様々なテクニックを用いてデータ量を擬似的に増やす方法である．
MDA-ASAS（Multiple Data-Automatic Short Answer Scoring）\cite{asag10}は，3つのデータ拡張手法を提案している．
具体的には，言語処理タスクで広く利用されている「逆翻訳」（例えば，文章を英語から日本語に翻訳し，それを英語に翻訳し直すことで言い回しの異なる同等の意味の文を作成する方法）を中心的に使用して，元のデータセットを8倍以上の量に拡張する手法を提案している．

\subsection*{小論文などの長めの回答文を対象とした深層学習自動採点モデル}\label{sec:other_aes}

\subsubsection*{階層型畳み込みニューラルネットワークモデル}

Dong and Zhang\cite{HierarchicalModel}は，文章の階層構造をモデル化する手法を提案した．
具体的には，単語レベルの畳み込みニューラルネットワーク（CNN: Convolutional Neural Network）に文単位で文章を入力し，文単位の分散表現ベクトルを作成する．
そして，文単位の分散表現ベクトルの系列を別のCNNに入力することで，文章全体の分散表現ベクトルを作成し，それを用いて得点を予測する手法である．

\subsubsection*{SKIPFLOWモデル}

文章の質を表す指標の一つである「一貫性」に着目したモデルとして，Tayら\cite{SKIPFLOWmodel}は，文章中の離れた位置の意味的一貫性を学習するSKIPFLOWと呼ばれる機構を組み込んだ自動採点モデルを提案している．具体的には， LSTMベースのモデルを自動採点モデルをベースに，ニューラルテンソルというアーキテクチャが導入されている．ニューラルテンソルは，LSTMの任意の異なる2時点の出力について類似度を計算する処理を行う．この機構によって，離れた位置の文脈の類似性を捉えることを目指している．

\subsubsection*{ハイブリッド手法}

2章で述べたように，自動採点モデルは，特徴量ベースの手法と深層学習ベースの手法に大別されるが，近年では，それらを統合したハイブリッド手法も提案されている．
例えば，Dasguptaら\cite{Hybridmodel1}が提案するモデルは，回答文の単語系列を処理するLSTMと，人手で設計された文レベルの特徴量系列を処理するLSTMを統合したモデルとして定式化されている．
また，Utoら\cite{Hybridmodel2}は，より単純なハイブリッド手法として，深層学習ベースの自動採点モデルで得られる文章レベルの分散表現ベクトルに，人手で設計した文章レベルの特徴量を連結し，それを用いて最終的な得点を予測するモデルを提案している．

\subsubsection*{項目反応理論を用いた複数自動採点モデルの統合手法}

Aomiら\cite{asag6}\cite{aomi2}は，複数の自動採点モデルの予測値を項目反応理論を用いて統合する手法を提案している．
特徴の異なる複数の自動採点モデルを，それぞれの採点特性を考慮して統合することで，単一の自動採点モデルを利用する場合に比べて，高精度な自動採点を実現している．

\subsubsection*{評価者バイアスに頑健な自動採点手法}

大規模な記述式試験では，多数の評価者が分担して採点を行うことが一般的であり，その場合，個々の答案に与えられる得点には評価者の特性（甘さ/厳しさなど）に依存したバイアスが反映されてしまう．
バイアスの影響を受けた訓練データを用いた場合，訓練された自動採点モデルにもそのバイアスの影響が引き継がれてしまい，自動採点の性能が低下する．
この問題を解決するために，Uto \& Okano\cite{asag8,okano,okano2}は，評価者特性の影響を取り除いて真の得点を推定できる項目反応モデル\cite{okano3,okano4}を利用した自動採点手法を提案している．
この手法により，様々な自動採点モデルにおいて，訓練データを作成する際の評価者の割り当て方法によらない，安定的な自動採点モデルの構築が実現されている．

%
%%% 謝辞
\chapter*{謝辞}
本研究を行うにあたり，親切にご指導を頂いた宇都雅輝准教授，植野真臣教授に厚く御礼を申し上げます．ゼミ発表中のコメントが大変励みになりました．また，プログラミングまわりの相談に際して岡野さん，柴田さんには重要な助言を賜りました．お陰で問題点を修正することができました．研究室の皆様に深く感謝申し上げます．
%
%
%%% 参考文献
\bibliographystyle{sieicej}
\bibliography{myrefs}
%
\newpage
\printindex
%
\end{document}
